\begin{center}
\begin{minipage}{0.92\textwidth}
\small
\noindent\textbf{Abstract}\quad
Contemporary computational neuroscience and deep learning have evolved along two largely disconnected trajectories. Macroscopic deep-learning architectures such as Transformers employ non-Markovian self-attention to model long-range temporal dependencies, yet operate with few biophysical constraints and high computational costs. Microscopic spiking-neuron models such as the Leaky Integrate-and-Fire (LIF) neuron respect physical principles and metabolic efficiency, but their strictly Markovian dynamics limit the extraction of higher-order temporal patterns. This paper proposes the \textbf{Temporal Attention Neuron (TAN)}, a single-neuron model that combines a Boltzmann-weighted temporal attention kernel with neuromodulatory surprise gating on top of a classical leaky integrator. TAN retains the local, Markovian physical substrate of LIF while adding a non-Markovian temporal receptive field driven by a deviation-from-history signal. Our experiments show that habituation, sparse sensory gating under noise, and adaptive niche localization emerge from the dynamics of this system, without any hard-coded adaptation module or pre-programmed navigation algorithm.

The paper presents a complete empirical and theoretical analysis organised in three layers. \textit{Emergent dynamics} (single-neuron and single-agent levels) demonstrate four behavioural phenomena: spontaneous habituation under sustained stimulation, sparse sensory gating under extreme noise, niche capture in a clean gradient, and hovering in a noisy environment that is phenomenologically analogous to \textit{E. coli} chemotaxis. \textit{Experimental evaluation} provides systematic validation across 50 random seeds: TAN achieves lower error than LIF on the two stochastic tasks with very large effect sizes ($p<10^{-30}$, Cohen's $|d|>5$); the two deterministic tasks produce identical trajectories across seeds, so no statistical test is applicable there, but TAN still achieves lower final error. TAN matches trained RNN and GRU baselines on three of four tasks \textit{without any training}; ablation results are consistent with a functional specialization of the three internal modules (surprise, attention, gate); and parameter sweeps show relatively broad performance plateaus over the tested ranges of TAN's unique parameters ($\lambda$, $\beta$). \textit{Dynamical and information-theoretic analysis} provides a mechanism for the observed behaviour: a fixed-point proof shows that the zero-membrane state is an attracting fixed point under sustained constant input; a 2D state-space projection reveals that trajectories converge to the origin; and a three-layer information cascade $I(z; S) \to I(z; A) \to I(z; y)$ shows that mutual information between noise and the spike output is lower than between noise and the internal current, consistent with information loss introduced by the downstream nonlinear gating and spike-generation stages. Together, the results suggest that QKV-like temporal weighting and deviation-driven gating can be combined at the single-neuron level, providing a computational primitive that may be potentially advantageous for neuromorphic implementations.

\vspace{0.6em}
\noindent\textbf{Keywords}\quad Spiking neural networks; Temporal attention; Surprise gating; Habituation; Neuromorphic computing; Artificial life
\end{minipage}
\end{center}
